% !PRINTED 654
\weberpage{654}
\begin{center}
\S\ 184. \gesperrt{Basis Form and Ramification Functional.}
\end{center}

By the basis form of the field $\Omega$ we shall understand the linear
 form:
\weq{1}{
\tau = \omega_1t_1 + \omega_2t_2 + \cdots + \omega_nt_n\,,
}
in which $t_1, t_2, \ldots, t_n$ are the functional variables, and whose coefficients
 $\omega_1, \omega_2, \ldots, \omega_n$ form a minimal basis of $\Omega$. This
form is the greatest common divisor of $\omega_1, \omega_2, \ldots, \omega_n$
and consequently a \gesperrt{unit.} From $\tau$ one can derive all integral
functions in $\Omega$ by putting for the variables $t_1,
t_2, \ldots, t_n$ integral functions in $Z$.

If we put
\weq{2}{
F(t) = N(t - \tau) = t^n + A_1t^{n-1} + \cdots + A_{n-1}t + A_n\,,
}
then the $A_1, A_2, \ldots, A_n$ are holomorphic integral functionals in
$\overline{Z}$, and $\tau$ satisfies the equation of degree $n$:
\weq{3}{
F(\tau) = 0.
}

If $\pi$ is a prime functional and
\[
N_a(\pi) = z - c\,,
\]
then by \S\ 182, 6. one can determine the \gesperrt{constants} $c_1, c_2, \ldots, c_n$ from
the congruences
\weq{4}{
\omega_1 \equiv c_1, \ \omega_2 \equiv c_2, \ \ldots, \ \omega_n \equiv c_n \quad (\mathrm{mod}\ \pi)
}
and if one puts
\weq{5}{\tau_0 = c_1t_1 + c_2t_2 + \cdots + c_nt_n}
then $\tau_0$ is a \gesperrt{constant functional}, and $\tau - \tau_0$ is
divisible by $\pi$. It is now to be proved:

3. \gesperrt{The prime functional $\pi$ is the greatest common
 divisor of $\tau - \tau_0$ and $z - c$:}
\weq{6}{
\pi = u(\tau - \tau_0) + v(z - c)\,,
}
\gesperrt{and if therefore $(z - c)$ is divisible by a higher than the
first power of $\pi$, then $\tau - \tau_0$ is divisible only by
the first power of $\pi$.}

First of all it follows from the definition \S\ 181 that
\weq{7}{
\tau - \tau_0 = t_1(\omega_1 - c_1) + t_2(\omega_2 - c_2) + \cdots + t_n(\omega_n - c_n)
}
is the greatest common divisor of $(\omega_1 - c_1), (\omega_2 - c_2), \ldots,
(\omega_n - c_n)$. If now
\weq{8}{
\omega = x_1\omega_1 + x_2\omega_2 + \cdots x_n\omega_n \equiv 0 \quad (\mathrm{mod}\ \pi)\,,
% [sic] printed "x_2\omega_2 + \cdots x_n\omega_n" with no "+" after \cdots, the sum requires "+ \cdots + x_n\omega_n" as in (1), (5) and (7)
}
%CONT>>
